\documentclass[11pt,a4paper]{article}
\usepackage[margin=1in]{geometry}
\usepackage{hyperref}
\usepackage{enumitem}
\usepackage{graphicx}
\usepackage{xcolor}



% Variable: Lab Instructor
\makeatletter \newcommand{\BvrSetLabInstructor}[1]{%
  \gdef\Bvr@LabInstructor{#1}}
% default
\newcommand{\Bvr@LabInstructor}{Jhonsy Bansal}
% public accessor
\newcommand{\BvrLabInstructorValue}{\Bvr@LabInstructor}
\makeatother

% Add subtitle
\usepackage{titling}
% -- Set title bold
\pretitle{\begin{center}\bfseries\fontsize{18bp}{18bp}\selectfont}
\makeatletter
\newcommand{\BvrSubtitle}[1]{%
  \posttitle{%
    \par\end{center}
    \begin{center}\large#1\end{center}
    \vskip 0.5em}%
}
\makeatother

% Hints in the section title(s)
\newcommand{\BvrHint}[1]{%
  {\normalfont\normalsize\itshape\hfill #1}}

% -----------------------------------------------------

\title{Camera-Based Weightlifting/Exercise Form Analysis and Feedback System}

\BvrSetLabInstructor{Jhonsy Bansal}

\BvrSubtitle{%
  Project Proposal for UCS503  \\
  Submitted to: \BvrLabInstructorValue \\ \bfseries
  Thapar Institute of Engineering and Technology%
}

\author{
  Lakshay Gambhir \; (1024031026) \\
  Gagandeep Singh \; (1024031027) \\
  Divjot Singh \; (1024031028)
}

\date{\today}

\begin{document}
\maketitle

\section{Higher-order goal%
  \BvrHint{\ldots secondary framing}}
This project aims to improve individual training safety and consistency by
giving fitness practitioners automated, camera-based feedback on their
weightlifting technique. While the broader motivation is accessible,
coach-free technique guidance, the immediate engineering objective is to
translate \emph{pose estimation and geometric biomechanical analysis} into a
measurable, deployable software solution that runs on hardware most users
already own -- a laptop webcam.

\section{Time-to-value %
\BvrHint{\ldots primary heuristic}}
\begin{itemize}[leftmargin=0.7cm]
    \item We will deliver a working, testable increment quickly by building
    one exercise (Squat) fully end-to-end -- pose estimation, rep detection,
    form rules, and feedback -- before extending to Deadlift and Bench Press.
    \item We will prioritize a verifiable feedback loop over exhaustive
    exercise coverage, since a single reliable exercise pipeline validates
    the full architecture and de-risks the remaining two.
    \item Success in the first iteration is defined as: correct rep counting
    and at least one correctly detected form violation on self-recorded
    Squat footage, measured against manual review.
\end{itemize}

\section{Problem Statement}
Recreational and intermediate lifters who train without a coach or training
partner have no objective, repeatable way to assess their own exercise
technique. Common pain points include:
\begin{itemize}[leftmargin=0.7cm]
    \item \textbf{No real-time awareness:} a lifter cannot see their own
    knee alignment, depth, or back angle while performing a lift.
    \item \textbf{Manual video review is slow:} filming and reviewing sets
    after the fact is time-consuming and requires technique knowledge to
    interpret correctly.
    \item \textbf{No structured history:} without dedicated tooling, there
    is no consistent record of technique quality across sessions, so
    progress (or regression) goes unnoticed.
    \item \textbf{Hardware barrier:} wearable-sensor or motion-capture-based
    solutions require equipment beyond what a typical home or gym user owns.
\end{itemize}

\textbf{Why this matters now (contextual relevance):} lightweight, real-time
pose-estimation models (e.g., MediaPipe Pose) now run efficiently on
consumer laptops without specialized hardware, making a camera-only
technique-feedback tool practically buildable within a single-semester
timeframe.

\section{Proposed Solution %
\BvrHint{\ldots Software Proposition}}
\subsection{Overview}
Build a \textbf{browser-based} ``Form Capture-to-Feedback'' system with the
following components:
\begin{itemize}[leftmargin=0.7cm]
    \item \textbf{Exercise selection \& camera capture:} user selects
    Squat, Deadlift, or Bench Press and grants webcam access.
    \item \textbf{Real-time pose pipeline:} pose estimation, landmark
    smoothing, and joint-angle calculation running client-side.
    \item \textbf{Rep detection \& form analysis:} a per-exercise movement
    state machine counts repetitions; a geometric rule engine flags common
    technique issues per rep.
    \item \textbf{Feedback \& scoring:} prioritized real-time text feedback
    during the set, plus a post-session form score and summary.
    \item \textbf{History dashboard:} session-level metrics (not video) are
    stored and displayed as trends over time.
\end{itemize}

\subsection{Core workflow}
\begin{enumerate}[leftmargin=0.7cm]
    \item User selects an exercise and starts a session after a camera
    visibility check.
    \item System tracks pose per frame, verifies the movement matches the
    selected exercise, and advances a rep-detection state machine.
    \item On each completed rep, form rules evaluate the rep's angle
    history and the top-priority issue (if any) is shown to the user.
    \item On session end, aggregate metrics and a form score are computed
    and saved; the dashboard is updated.
\end{enumerate}

\subsection{Operational constraints for an educational, single-semester setting}
\begin{itemize}[leftmargin=0.7cm]
    \item Keep the solution usable on a standard laptop webcam; no
    specialized sensors, depth cameras, or GPUs required.
    \item Avoid dependence on a large proprietary or novel ML training
    pipeline; focus on pose landmarks, geometry, and rule correctness.
    \item Design for deployability using common web hosting and a managed
    relational database.
\end{itemize}

\section{Solution Approach %
\BvrHint{\ldots Engineering focus}}
This is an engineering course project, so we focus on a correct,
modular real-time pipeline rather than novel computer-vision research.

\subsection{Form analysis strategy %
\BvrHint{\ldots non-research}}
Form analysis will be rule-based, not classifier-based:
\begin{itemize}[leftmargin=0.7cm]
    \item Extract exercise-specific joint angles (e.g., knee, hip, torso for
    Squat) from smoothed pose landmarks.
    \item Compare each rep's angle trajectory against configurable
    thresholds (e.g., minimum depth, maximum knee drift) without claiming a
    validated clinical or biomechanical model.
    \item Surface only the single highest-priority violation per rep to the
    user, never an unranked list of warnings.
\end{itemize}

\subsection{Web architecture %
\BvrHint{\ldots browser-based solution}}
A client-heavy three-tier architecture:
\begin{itemize}[leftmargin=0.7cm]
    \item \textbf{Frontend (client):} React UI running camera capture,
    MediaPipe Pose, rep detection, form analysis, and feedback in real time.
    \item \textbf{Backend API:} FastAPI service for authentication, session
    submission, and history/dashboard queries.
    \item \textbf{Database:} PostgreSQL storing structured session, rep,
    and form-issue records -- never workout video.
\end{itemize}

\subsection{CI/CD and fast delivery enabler}
A CI/CD pipeline will be used to:
\begin{itemize}[leftmargin=0.7cm]
    \item Automatically run tests (angle calculation, state-machine logic)
    on every change.
    \item Produce repeatable deployment artifacts for the frontend and
    backend.
    \item Enable safe, frequent releases as each exercise (Squat, then
    Deadlift, then Bench Press) is completed.
\end{itemize}

\section{Evaluation Criterion %
\BvrHint{\ldots measurable and attributable}}
We define success using measurable metrics tied to the core pipeline.

\subsection{Primary evaluation metric}
\textbf{Rep-counting accuracy:} agreement between system-counted repetitions
and manually counted repetitions on self-recorded test sets.
\begin{itemize}[leftmargin=0.7cm]
    \item Target: high agreement across a self-recorded test set covering
    all three exercises under normal camera conditions.
    \item Attribution: measured by comparing logged rep counts against a
    human reviewer's manual count on the same recordings.
\end{itemize}

\subsection{Secondary evaluation metrics}
\begin{itemize}[leftmargin=0.7cm]
    \item \textbf{Form-detection agreement:} percentage of deliberately
    introduced technique errors (e.g., shallow squat, rounded back) that
    are correctly flagged by the rule engine.
    \item \textbf{Visibility-handling correctness:} whether the system
    correctly pauses/warns rather than emits confident feedback under poor
    lighting or occlusion.
    \item \textbf{Session completeness:} percentage of sessions that save
    successfully to the database without data loss.
    \item \textbf{Reliability:} basic uptime/availability of the
    staging/production deployment during testing windows.
\end{itemize}

\subsection{Pilot validation plan %
\BvrHint{\ldots high-level}}
\begin{itemize}[leftmargin=0.7cm]
    \item Record a self-collected test set (self and, where possible, a
    small number of additional volunteers) performing clean and
    deliberately flawed reps for each exercise.
    \item Compare system output (rep counts, flagged issues) against
    manual review of the same footage.
    \item Where accessible, cross-check Squat detections against a public
    labelled dataset (e.g., Fitness-AQA) as an external reference point.
\end{itemize}

\section{Scalability %
\BvrHint{\ldots with foresight}}
The proposition should scale in both theoretical and practical senses.

\begin{enumerate}[leftmargin=0.7cm]
    \item \textbf{Scalable (theoretically):} the system should support
    growth in users and sessions by:
    \begin{itemize}[leftmargin=0.7cm]
        \item keeping real-time pose processing client-side, so backend
        load scales only with session-submission and dashboard traffic,
        \item decoupling frontend and backend,
        \item using indexed queries for session/history lookups.
    \end{itemize}

    \item \textbf{Operationally deployable (practically):} the system
    should run on common web hosting:
    \begin{itemize}[leftmargin=0.7cm]
        \item managed PostgreSQL database,
        \item standard frontend hosting for the React app,
        \item CI/CD pipeline for staging/production releases.
    \end{itemize}
\end{enumerate}

\section{Engine availability heuristic}
A mature set of ``engines'' (tooling/libraries) is available to implement
this as a browser-based project:
\begin{itemize}[leftmargin=0.7cm]
    \item MediaPipe Pose (Tasks API, JavaScript) for real-time, on-device
    pose landmark extraction.
    \item Web framework for REST APIs (FastAPI).
    \item Managed relational database (PostgreSQL).
    \item Authentication approach (token-based).
    \item Charting library for the React dashboard.
    \item Test framework for unit/integration tests (state machine, angle
    math, API endpoints).
    \item CI runner and staging deployment target.
\end{itemize}
We will use these mature, existing components to focus on the pipeline
design, rule correctness, and measurement, rather than inventing new
computer-vision infrastructure.

\section{Project scope and deliverables %
\BvrHint{\ldots iteration-friendly}}
\subsection{Initial deliverable %
\BvrHint{\ldots first iteration}}
\begin{itemize}[leftmargin=0.7cm]
    \item Camera capture + pose landmark extraction with live overlay
    \item Squat rep-detection state machine with smoothing/hysteresis
    \item Squat form rules (depth, knee alignment, torso angle) with
    prioritized real-time feedback
    \item Session save to backend (rep count, metrics, issues -- no video)
    \item CI: build + unit tests for angle/state-machine logic
\end{itemize}

\subsection{Subsequent deliverables}
\begin{itemize}[leftmargin=0.7cm]
    \item Deadlift and Bench Press state machines and form rules
    \item Configurable weighted form score
    \item History dashboard with session trends and common-issue summaries
    \item Calibration against a self-recorded test set (and, where
    accessible, an external dataset reference check)
    \item CD to staging with a repeatable deployment pipeline
\end{itemize}

\section{Risks and mitigations}
\begin{itemize}[leftmargin=0.7cm]
    \item \textbf{Single-camera accuracy limits:} depth-dependent metrics
    (e.g., precise bar path) cannot be reliably measured from one webcam;
    such metrics will be explicitly labelled as approximate or unavailable
    rather than presented as accurate.
    \item \textbf{Noisy pose data / false rep counts:} mitigated via
    landmark smoothing and hysteresis thresholds in the state machine.
    \item \textbf{No large labelled weightlifting-form dataset:} V1 does
    not depend on one -- form analysis is rule-based; a self-recorded test
    set is used for calibration and evaluation instead.
    \item \textbf{Scope creep beyond three exercises:} mitigated by
    building and validating Squat completely before starting Deadlift or
    Bench Press (see Section 2).
    \item \textbf{Privacy concerns:} mitigated by never storing workout
    video -- only structured, derived metrics are persisted.
\end{itemize}

\section{Summary}
This project proposes a deployable, browser-based system that gives
weightlifters automated, real-time technique feedback for Squat, Deadlift,
and Bench Press using only a standard laptop webcam. It meets the course
criteria by emphasizing time-to-value through an incremental,
one-exercise-at-a-time build order, using measurable evaluation criteria
(rep-counting accuracy and form-detection agreement against manual review),
and remaining focused on engineering correctness and modular pipeline
design rather than novel computer-vision research or dependence on a
large proprietary dataset.

\end{document}
